\chapter{Incremental Search: \lpastar and \dstarlite}
\label{ch:ch05}
% Function names for pseudocode (unique across the book: prefixed Lpa/Dsl/Inc).
\SetKwFunction{LpaCalculateKey}{CalculateKey}
\SetKwFunction{LpaInitialize}{Initialize}
\SetKwFunction{LpaUpdateVertex}{UpdateVertex}
\SetKwFunction{LpaComputeShortestPath}{ComputeShortestPath}
\SetKwFunction{LpaMain}{Main}
\SetKwFunction{DslCalculateKey}{CalculateKey}
\SetKwFunction{DslInitialize}{Initialize}
\SetKwFunction{DslUpdateVertex}{UpdateVertex}
\SetKwFunction{DslComputeShortestPath}{ComputeShortestPath}
\SetKwFunction{DslMain}{Main}
\SetKwFunction{IncTopKey}{TopKey}
\SetKwFunction{IncPop}{Pop}
\SetKwFunction{IncInsert}{Insert}
\SetKwFunction{IncRemove}{Remove}

\begin{objectives}
\begin{itemize}
  \item Explain why re-running \astar after a small map change wastes work, and what an
    \emph{incremental} search reuses instead.
  \item Define the $\gcost$- and $\rhs$-values of \lpastar, local consistency, over- and
    underconsistency, and the two-component priority key, and trace \lpastar by hand on a
    small grid.
  \item Explain why \dstarlite searches backward from the goal, what the key modifier
    $k_m$ repairs when the robot moves, and reproduce its pseudocode.
  \item State the conditions under which \dstarlite returns the same path costs as
    re-running \astar, and bound its number of vertex expansions.
  \item Implement \dstarlite in \python, change obstacles while the agent executes its
    plan, and measure the saving relative to re-running \astar.
\end{itemize}
\end{objectives}

% ---------------------------------------------------------------------
\section{Why this matters for a drone swarm}
\label{sec:ch05-motivation}

Picture one drone of a swarm on its nominal route through a warehouse. The route was
computed by the global planner (\cref{ch:ch09,ch:ch10}) before take-off, on a grid of
$200 \times 200$ cells at ten flight levels. Halfway through the mission the prediction
layer (\cref{ch:ch18,ch:ch20}) announces that an unknown drone will cross the corridor
ahead within the next few seconds. The local layer (\cref{ch:ch13,ch:ch14}) can dodge for
a moment, but the intruder is going to hover in the corridor, and the local manoeuvre
cannot return the drone to its nominal route. Now the replanning layer must deliver a new
route from the drone's \emph{current} position, and it has one control cycle,
say $50$~ms, to do so.

Re-running \astar (\cref{ch:ch04}) from scratch is the obvious answer. It is also
wasteful: of the $400\,000$ cells of the map, only a few dozen have changed. Every
distance that \astar computed last time is still valid, except in a small region around
the intruder. An \emph{incremental} search keeps the old search results and repairs only
what the change invalidated. \dstarlite, the algorithm of this chapter, does exactly
that, and on top of it handles the fact that the drone keeps moving while it replans.
It is the replanning algorithm of the hybrid architecture in \cref{ch:ch24}, and the
training plan rates it \priority{Essential}; its conceptual parent \lpastar is rated
\priority{Medium}.

The chapter proceeds in two halves. The first half develops \lpastar, which repairs a
shortest path from a fixed start to a fixed goal when edge costs change. It introduces
the two numbers stored per vertex ($\gcost$ and $\rhs$), the notion of local consistency,
and the two-component key that orders the work. The second half turns \lpastar into
\dstarlite by reversing the search direction and adding the key modifier $k_m$, so that
the start vertex may move. Both halves have a fully worked grid example whose numbers
come from the chapter's \python code, and the chapter ends with an experiment that
compares \dstarlite with repeated \astar on a $50 \times 50$ grid as obstacles appear
one by one.

% ---------------------------------------------------------------------
\section{The idea in plain words}
\label{sec:ch05-intuition}

\Cref{fig:ch05-idea} shows the situation of the motivation on a small grid. The drone
has a shortest path to its goal. An obstacle appears on that path. In the left panel,
\astar is re-run from scratch: it expands every cell shaded grey again, including
hundreds whose distance did not change at all. In the right panel, an incremental search
touches only the orange cells: the ones whose distance value became stale because the
obstacle removed the edge they relied on, plus the few cells needed to find the detour.

\begin{figure}[tb]
  \centering
  \inputfigure{ch05/idea}
  \caption{An obstacle (black, with a cross) appears on the planned path (blue). Left:
  \astar re-run from scratch expands all grey cells again. Right: \dstarlite expands only
  the orange cells, those whose goal distance was invalidated or is needed for the
  detour (dashed). The numbers of expansions are printed under each panel.}
  \label{fig:ch05-idea}
\end{figure}

How does the incremental search know which cells are stale? It stores \emph{two}
numbers per vertex instead of one. The first, $\gcost(s)$, is the distance value from
the last search, exactly as in \astar. The second, $\rhs(s)$, is what the neighbours of
$s$ currently say the value \emph{should} be: one step of lookahead, ``the best
neighbour's $\gcost$ plus the edge cost''. As long as the two numbers agree, nothing at
$s$ needs attention. When an edge cost changes, the $\rhs$-value of the affected vertex
changes but its $\gcost$-value does not, and the disagreement flags the vertex as
\emph{locally inconsistent}. Inconsistent vertices go into a priority queue, ordered by
the same kind of key that \astar uses ($\gcost + \hcost$), and are processed in that
order. Processing a vertex makes it consistent and may make its neighbours inconsistent,
so the repair spreads outward from the change, but only as far as distance values
actually differ, and only until the goal is consistent and no queued key is smaller than
the goal's key. Cells far from the change, or cells whose distance happens to stay the
same, are never touched.

\begin{keyidea}
Store, for every vertex, the value from the last search ($\gcost$) next to the value its
neighbours currently imply ($\rhs$). Only vertices where the two disagree need work.
Process them in \astar order, from the change outward, and stop as soon as the query
vertex agrees with its neighbours and nothing with a smaller key is left. The result
equals what \astar would compute from scratch, at a fraction of the expansions when the
change is small.
\end{keyidea}

\dstarlite adds one more idea. A moving robot changes the \emph{start} of the search at
every step. If distance values were measured from the start, every one of them would
change with every step. Distances \emph{to the goal} do not change when the robot moves.
So \dstarlite searches backward from the goal, and the robot's motion only affects the
heuristic part of the keys, which a single running offset $k_m$ can absorb.

% ---------------------------------------------------------------------
\section{Problem statement and notation}
\label{sec:ch05-problem}

We use the graph notation of \cref{ch:ch02,ch:ch03}: a finite directed graph $G=(V,E)$
with edge costs $c(u,v) \in (0,\infty]$ for $(u,v)\in E$, a start vertex
$s_{\mathrm{start}}$ and a goal vertex $s_{\mathrm{goal}}$. $\Pred(s)$ and $\Succ(s)$
are the predecessors and successors of $s$. A cost of $\infty$ means that the edge is
currently unusable; we write $c(u,v)=\infty$ also for pairs that are not adjacent, which
keeps formulas uniform. The true distance from $u$ to $v$ under the current costs is
$\dist(u,v)$, and $\dist(u,v)=\infty$ if no path exists.

\begin{definition}[Incremental shortest-path problem]
\label{def:ch05-problem}
An \textbf{incremental shortest-path problem}\index{incremental search}\index{shortest
path!incremental} is a sequence of shortest-path queries between the same two vertices
$s_{\mathrm{start}}$ and $s_{\mathrm{goal}}$ on a graph whose edge costs change between
queries. An \textbf{incremental search algorithm} answers each query using information
saved from the previous queries, and it must return a path of cost
$\dist(s_{\mathrm{start}}, s_{\mathrm{goal}})$ under the costs current at the time of
the query.
\end{definition}

The grids of this chapter are 4-connected with unit costs (\cref{ch:ch02}). A cell is
either \emph{free} or \emph{blocked}. Blocking a cell $b$ is modelled as an edge-cost
change: \emph{every} edge that enters $b$ and every edge that leaves $b$ changes from
cost $1$ to cost $\infty$. Unblocking reverses the change.\index{obstacle!as infinite
edge cost} We keep blocked cells as vertices of the graph, with all their edges at
$\infty$, rather than deleting them; that way the graph structure is fixed and only the
costs vary, which is exactly what the algorithms assume. The heuristic on these grids is
the Manhattan distance, which is consistent (\cref{ch:ch04}).

\lpastar stores two values per vertex.

\begin{definition}[$\gcost$-value and $\rhs$-value]
\label{def:ch05-rhs}
For every vertex $s\in V$, $\gcost(s)\in[0,\infty]$ is the \textbf{g-value}\index{LPA*!g-value},
an estimate of $\dist(s_{\mathrm{start}},s)$ that is only changed when $s$ is expanded.
The \textbf{rhs-value}\index{LPA*!rhs-value}\index{rhs-value} (``right-hand side'',
after the Bellman equation it comes from) is the one-step lookahead
\begin{equation}
\label{eq:ch05-rhs}
\rhs(s) =
\begin{cases}
0 & \text{if } s = s_{\mathrm{start}},\\[2pt]
\displaystyle\min_{s'\in\Pred(s)} \bigl(\gcost(s') + c(s',s)\bigr) & \text{otherwise.}
\end{cases}
\end{equation}
\end{definition}

The $\rhs$-value is always kept up to date with respect to the current $\gcost$-values
of the predecessors and the current edge costs; the $\gcost$-value is allowed to lag
behind. The gap between them is the central concept.

\begin{definition}[Local consistency]
\label{def:ch05-consistency}
A vertex $s$ is \textbf{locally consistent}\index{local consistency}\index{LPA*!local
consistency} if $\gcost(s) = \rhs(s)$, and \textbf{locally inconsistent} otherwise. An
inconsistent vertex is \textbf{overconsistent}\index{overconsistent vertex} if
$\gcost(s) > \rhs(s)$ (its neighbours now offer a cheaper way to reach it than it
knows) and \textbf{underconsistent}\index{underconsistent vertex} if
$\gcost(s) < \rhs(s)$ (it believes in a cost that its neighbours can no longer
support).
\end{definition}

A freshly discovered vertex has $\gcost=\infty$ and a finite $\rhs$, so it is
overconsistent; this is the only kind of inconsistency \astar ever meets. Underconsistent
vertices appear when a cost \emph{increases}, for example when an obstacle appears on
the path: the vertex behind the obstacle still holds its old, now too optimistic
$\gcost$-value. The following proposition says why consistency is the right target.

\begin{proposition}[Consistency everywhere means correct distances]
\label{thm:ch05-consistent-correct}
Suppose all edge costs are positive. If every vertex is locally consistent, then
$\gcost(s) = \dist(s_{\mathrm{start}}, s)$ for all $s \in V$, and a shortest path from
$s_{\mathrm{start}}$ to any reachable $s$ is obtained by moving backward from $s$ to a
predecessor $s'$ that attains the minimum in \cref{eq:ch05-rhs}, until
$s_{\mathrm{start}}$ is reached.
\end{proposition}
\begin{proof}
If all vertices are consistent, the $\gcost$-values satisfy $\gcost(s_{\mathrm{start}})=0$
and $\gcost(s)=\min_{s'\in\Pred(s)}(\gcost(s')+c(s',s))$ for all other $s$: the Bellman
equations of the shortest-path problem. With positive edge costs, the true distances are
the unique solution of these equations (\cref{ch:ch03}), so $\gcost \equiv \dist(s_{\mathrm{start}},\cdot)$.
Following a minimising predecessor decreases $\gcost$ by exactly the edge cost at every
step, so the walk terminates at $s_{\mathrm{start}}$ and has cost $\gcost(s)$.
\end{proof}

\lpastar does not make \emph{all} vertices consistent; like \astar it stops early. What
it guarantees is that the goal is consistent and that no queued vertex could still
improve the goal, and that suffices for the goal's $\gcost$-value and the path traced
from it to be correct (\cref{thm:ch05-correct}). To decide the order of work it uses a
key with two components.

\begin{definition}[Priority key]
\label{def:ch05-key}
The \textbf{key}\index{key!two-component}\index{LPA*!key} of a vertex $s$ is the pair
\begin{equation}
\label{eq:ch05-key}
k(s) = \bigl[\,k_1(s);\; k_2(s)\,\bigr] =
\bigl[\,\min(\gcost(s),\rhs(s)) + \hcost(s);\;\; \min(\gcost(s),\rhs(s))\,\bigr],
\end{equation}
where $\hcost(s)$ estimates $\dist(s,s_{\mathrm{goal}})$ and is nonnegative and
consistent: $\hcost(s_{\mathrm{goal}})=0$ and $\hcost(s)\le c(s,s')+\hcost(s')$ for all
$(s,s')\in E$. Keys are compared \textbf{lexicographically}: $k \le k'$ iff $k_1 < k'_1$,
or $k_1 = k'_1$ and $k_2 \le k'_2$.
\end{definition}

Read $k_1$ as the $\fcost$-value of \astar and $k_2$ as its $\gcost$-value used for
tie-breaking, both computed from the \emph{smaller} of the two stored values. For an
overconsistent vertex the smaller value is $\rhs$, the tentative new distance, exactly
as in \astar. For an underconsistent vertex it is the old $\gcost$, which is too small.
Using the too-small value on purpose puts the vertex \emph{early} in the queue, before
every vertex whose value depends on it, so that the stale value is retracted before it
can be propagated further. \Cref{fig:ch05-consistency} illustrates the three states.

\begin{figure}[tb]
  \centering
  \inputfigure{ch05/consistency}
  \caption{Local consistency on a three-vertex fragment (edge costs on the arrows,
  $\gcost/\rhs$ inside the vertices). Left: all consistent. Middle: the edge into $v$
  became expensive, so $\rhs(v)$ rose and $v$ is underconsistent; its key uses the old
  $\gcost(v)=3$ and puts $v$ before $w$, whose $\rhs$ still relies on $\gcost(v)$. Right:
  a new cheap edge into $v$ makes $v$ overconsistent; its key uses $\rhs(v)=1$.}
  \label{fig:ch05-consistency}
\end{figure}

The priority queue $U$ holds exactly the locally inconsistent vertices, each with the
key it had when it was last inserted. $U.\IncTopKey{}$ returns the smallest key in $U$
(or $[\infty;\infty]$ if $U$ is empty), $U.\IncPop{}$ removes and returns the vertex
with that key, $U.\IncInsert{s, k}$ inserts $s$ with key $k$, and $U.\IncRemove{s}$
removes $s$.

% ---------------------------------------------------------------------
\section{The algorithm: \lpastar}
\label{sec:ch05-lpastar}

\Cref{alg:ch05-lpastar} is Lifelong Planning \astar (\lpastar) of Koenig, Likhachev and
Furcy \cite{koenig2004lpa}, in the form of the original paper. It searches
\emph{forward} from $s_{\mathrm{start}}$; $\gcost$-values are start distances, and
$\rhs$-values look at predecessors.\index{LPA*}\index{Lifelong Planning A*|see{LPA*}}

\begin{algorithm}[htb]
\caption{Lifelong Planning \astar (\lpastar) \cite{koenig2004lpa}}
\label{alg:ch05-lpastar}
\KwIn{graph $G$, costs $c$, start $s_{\mathrm{start}}$, goal $s_{\mathrm{goal}}$, consistent heuristic $\hcost$}
\KwOut{after every call of \LpaComputeShortestPath: $\gcost(s_{\mathrm{goal}})=\dist(s_{\mathrm{start}},s_{\mathrm{goal}})$ and a shortest path (\cref{thm:ch05-correct})}
\Fn{\LpaCalculateKey{$s$}}{
  \KwRet{$[\min(\gcost(s),\rhs(s)) + \hcost(s);\ \min(\gcost(s),\rhs(s))]$}\label{alg:ch05-lpastar:key}\;
}
\Fn{\LpaInitialize{}}{
  $U \gets \emptyset$\;
  \lForEach{$s \in V$}{$\rhs(s) \gets \gcost(s) \gets \infty$}
  $\rhs(s_{\mathrm{start}}) \gets 0$\label{alg:ch05-lpastar:init}\;
  $U.\IncInsert{$s_{\mathrm{start}}$, $[\hcost(s_{\mathrm{start}});\,0]$}$\;
}
\Fn{\LpaUpdateVertex{$u$}}{
  \lIf{$u \neq s_{\mathrm{start}}$}{$\rhs(u) \gets \min_{s' \in \Pred(u)} \bigl(\gcost(s') + c(s',u)\bigr)$\label{alg:ch05-lpastar:rhs}}
  \lIf{$u \in U$}{$U.\IncRemove{$u$}$\label{alg:ch05-lpastar:remove}}
  \lIf{$\gcost(u) \neq \rhs(u)$}{$U.\IncInsert{$u$, \LpaCalculateKey{$u$}}$\label{alg:ch05-lpastar:insert}}
}
\Fn{\LpaComputeShortestPath{}}{
  \While{$U.\IncTopKey{} < $ \LpaCalculateKey{$s_{\mathrm{goal}}$} \KwOr $\rhs(s_{\mathrm{goal}}) \neq \gcost(s_{\mathrm{goal}})$\label{alg:ch05-lpastar:while}}{
    $u \gets U.\IncPop{}$\label{alg:ch05-lpastar:pop}\;
    \uIf(\tcp*[f]{overconsistent}){$\gcost(u) > \rhs(u)$\label{alg:ch05-lpastar:over}}{
      $\gcost(u) \gets \rhs(u)$\;
      \lForEach{$s \in \Succ(u)$}{\LpaUpdateVertex{$s$}\label{alg:ch05-lpastar:oversucc}}
    }
    \Else(\tcp*[f]{underconsistent}){
      $\gcost(u) \gets \infty$\label{alg:ch05-lpastar:under}\;
      \lForEach{$s \in \Succ(u) \cup \set{u}$}{\LpaUpdateVertex{$s$}\label{alg:ch05-lpastar:undersucc}}
    }
  }
}
\Fn{\LpaMain{}}{
  \LpaInitialize{}\;
  \While{\KwTrue}{
    \LpaComputeShortestPath{}\label{alg:ch05-lpastar:compute}\;
    wait for changes in edge costs\;
    \ForEach{directed edge $(u,v)$ whose cost changed\label{alg:ch05-lpastar:scan}}{
      update the stored cost $c(u,v)$\;
      \LpaUpdateVertex{$v$}\label{alg:ch05-lpastar:changed}\;
    }
  }
}
\end{algorithm}

\paragraph{Walkthrough.}
\LpaInitialize sets every $\gcost$ and $\rhs$ to $\infty$ except
$\rhs(s_{\mathrm{start}})=0$ (line~\ref{alg:ch05-lpastar:init}); the start is thus the
only inconsistent vertex and the only entry of $U$, just as it is the only entry of
\Open at the start of \astar.

\LpaUpdateVertex{$u$} is the maintenance routine. Line~\ref{alg:ch05-lpastar:rhs}
recomputes $\rhs(u)$ from \cref{eq:ch05-rhs}; the start keeps $\rhs=0$ forever.
Lines~\ref{alg:ch05-lpastar:remove}--\ref{alg:ch05-lpastar:insert} then restore the
invariant ``$U$ contains exactly the inconsistent vertices, with current keys'': $u$ is
taken out, and put back with a fresh key only if it is inconsistent. Every path of the
algorithm that touches a $\gcost$-value or an edge cost calls \LpaUpdateVertex on the
vertices whose $\rhs$ could have changed, and nothing else ever modifies $\rhs$.

\LpaComputeShortestPath is the search proper. Its loop condition
(line~\ref{alg:ch05-lpastar:while}) mirrors the termination of \astar: continue while
some queued vertex has a key smaller than the goal's, or while the goal itself is
inconsistent. Each iteration pops the vertex $u$ with the smallest key
(line~\ref{alg:ch05-lpastar:pop}) and treats the two kinds of inconsistency
differently:
\begin{itemize}
  \item \emph{Overconsistent} ($\gcost(u) > \rhs(u)$, line~\ref{alg:ch05-lpastar:over}):
    accept the better value, $\gcost(u)\gets\rhs(u)$. This is the expansion of \astar.
    The successors of $u$ may now have a better predecessor, so they are updated
    (line~\ref{alg:ch05-lpastar:oversucc}). After this step $u$ is consistent.
  \item \emph{Underconsistent} ($\gcost(u) < \rhs(u)$, line~\ref{alg:ch05-lpastar:under}):
    the old value cannot be trusted, so it is thrown away, $\gcost(u)\gets\infty$.
    Now every successor that computed its $\rhs$ through $u$ must recompute it, and $u$
    itself must recompute its own $\rhs$ and re-enter the queue if it is not $\infty$
    (line~\ref{alg:ch05-lpastar:undersucc}, note the $\cup\set{u}$). After this step $u$
    is either consistent (both $\infty$) or overconsistent with a fresh, larger key, and
    will be expanded a second time later as an ordinary \astar expansion.
\end{itemize}

\LpaMain glues the pieces together: after each search, wait for the environment to
report changed edges, update their costs, and call \LpaUpdateVertex on the \emph{head}
$v$ of every changed edge (line~\ref{alg:ch05-lpastar:changed}), because it is $v$
whose $\rhs$ depends on $c(u,v)$. Then search again. The first call of
\LpaComputeShortestPath is exactly an \astar search (\cref{thm:ch05-expansions}); each
later call is a repair.

\paragraph{Why the two-component key works.}
Consider what the loop condition must guarantee: when it stops, the goal's
$\gcost$-value must equal the true distance. The first key component
$k_1 = \min(\gcost,\rhs) + \hcost$ is a lower bound on the cost of any path to the goal
through the queued vertex, so as long as $\IncTopKey{} \ge k(s_{\mathrm{goal}})$ no
queued vertex can lead to a cheaper path to the goal than the goal's own value, exactly
the argument that proves \astar optimal with a consistent heuristic (\cref{ch:ch04}).
The second component breaks ties in favour of \emph{smaller} $\min(\gcost,\rhs)$, which
matters more here than in \astar: among vertices with equal $k_1$, an underconsistent
vertex $u$ with a small old $\gcost(u)$ must be processed before a successor $w$ whose
$\rhs(w)=\gcost(u)+c(u,w)$ was derived from it. Because $c(u,w)>0$, $k_2(u) < k_2(w)$
whenever $k_1(u)=k_1(w)$, so the stale value is withdrawn first and $w$ is recomputed
from correct information. Without the second component, $w$ could be expanded with a
value that is about to be retracted, and would have to be expanded a third time.

\begin{pitfall}[Comparing keys as single numbers]
A key is a pair, and both components matter. Comparing only $k_1$ (or, worse, only
$k_2$) breaks the ordering argument above and can make \lpastar expand vertices with
stale values, or terminate with a wrong $\gcost(s_{\mathrm{goal}})$. In \python, store
keys as tuples \code{(k1, k2)}; tuple comparison is lexicographic, which is what the
algorithm needs. Also never compare keys for equality with floating-point tolerance
tricks: the algorithm only needs $<$ and $\le$.
\end{pitfall}

% ---------------------------------------------------------------------
\section{A worked example: \lpastar repairs a path}
\label{sec:ch05-lpastar-example}

\begin{example}[An obstacle appears on the path]
\label{ex:ch05-lpastar}
\Cref{fig:ch05-lpastar-example} shows a $5\times 3$ grid of free cells. Cells are named
$(x,y)$ with $x$ the column from the left and $y$ the row from the bottom, both starting
at $0$. The start is $S=(0,1)$, the goal is $T=(4,1)$, moves are 4-connected with unit
cost, and $\hcost$ is the Manhattan distance to $T$. Ties between equal keys are broken
toward the smaller $x$, then the smaller $y$. First \lpastar computes a shortest path.
Then cell $(3,1)$, which lies on the path, becomes blocked, and \lpastar repairs the
path.
\end{example}

\begin{figure}[tb]
  \centering
  \inputfigure{ch05/lpastar-example}
  \caption{\lpastar on \cref{ex:ch05-lpastar}. Each cell shows $\gcost/\rhs$; grey cells
  are locally consistent with finite values (expanded), blue cells are in the queue $U$,
  orange cells are underconsistent. Left: after the first search, path of cost $4$.
  Middle: right after cell $(3,1)$ is blocked and \LpaUpdateVertex has run on the
  affected vertices, before \LpaComputeShortestPath. Right: after the repair, the new path
  of cost $6$.}
  \label{fig:ch05-lpastar-example}
\end{figure}

\paragraph{The first search.}
\Cref{tab:ch05-lpastar-first} traces the first call of \LpaComputeShortestPath. Every
popped vertex is overconsistent, so this is an \astar search: the cells of the middle
row are expanded one after the other, each with key $[4;\,x]$, because for a cell on
the straight line $\gcost + \hcost = 4$. The side cells receive finite $\rhs$-values and
enter the queue with $k_1 = 6$, but they are never expanded: after $T$ is popped and made
consistent, the smallest queued key is $[6;1]$, which is not smaller than $k(T)=[4;4]$,
and the loop stops. Five expansions, ten cells left in the queue with correct
$\rhs$-values but $\gcost=\infty$.

\begin{table}[tb]
  \centering\small
  \caption{First call of \LpaComputeShortestPath in \cref{ex:ch05-lpastar}. ``Updated''
  lists the vertices on which \LpaUpdateVertex changed something, with their new
  $\rhs$-value and key.}
  \label{tab:ch05-lpastar-first}
  \begin{tabular}{@{}rllll@{}}
  \toprule
  Step & Popped $u$ (key) & Case & $\gcost(u)$ after & Updated: $\rhs$, key\\
  \midrule
  %%NUM-LPA-FIRST
  \bottomrule
  \end{tabular}
\end{table}

\paragraph{The change.}
Blocking $(3,1)$ changes eight directed edges: the four that enter $(3,1)$ and the four
that leave it. \LpaMain calls \LpaUpdateVertex on the head of each. For the entering
edges the head is $(3,1)$ itself: all its predecessors are now at cost $\infty$, so
$\rhs(3,1)=\infty$ while $\gcost(3,1)=3$; the cell is underconsistent and enters the
queue with key $[3+1;\,3]=[4;3]$. For the leaving edges the heads are the neighbours.
$T=(4,1)$ loses its only finite predecessor, $\rhs(T)=\infty$ against $\gcost(T)=4$:
underconsistent, key $[4;4]$. $(2,1)$ still has $S$'s side as predecessor and keeps
$\rhs=2$. $(3,0)$ and $(3,2)$ had $\rhs=4$ through $(3,1)$; now their only finite
predecessor is gone, $\rhs=\infty=\gcost$, and they \emph{leave} the queue. The middle
panel of \cref{fig:ch05-lpastar-example} shows this state: two underconsistent cells,
eight queued cells, and every $\gcost$-value untouched.

\paragraph{The repair.}
\Cref{tab:ch05-lpastar-repair} traces the second call. The two underconsistent cells
come first, because their keys use the old, small $\gcost$-values. Popping $(3,1)$ sets
$\gcost(3,1)=\infty$; its successors are updated but nothing changes, since their
$\rhs$-values already ignored $(3,1)$. Popping $T$ sets $\gcost(T)=\infty$, and now the
cells $(4,0)$ and $(4,2)$, whose $\rhs=5$ came through $T$, also drop out of the queue.
From here on the search is a plain \astar over the side rows: the queued cells are
expanded in order of their keys, all with $k_1=6$ and increasing $k_2$, the detour
$(2,0)\to(3,0)\to(4,0)$ is discovered, and $T$ is finally re-expanded with
$\gcost(T)=6$. The repair costs %%NUM-LPA-REPAIR expansions, of which the first two
are underconsistent. Compare this with \astar from scratch on the changed grid, which
expands %%NUM-ASTAR-AFTER cells: on such a small instance the saving is modest, because
almost every cell has $\fcost = 6$ and has to be looked at by any best-first search.
What \lpastar reuses are the $\gcost$-values of $S$, $(1,1)$ and $(2,1)$, which are
never expanded again. \Cref{sec:ch05-experiment} shows that on larger maps the reused
part dominates.

\begin{table}[tb]
  \centering\footnotesize
  \caption{Second call of \LpaComputeShortestPath in \cref{ex:ch05-lpastar}, after
  $(3,1)$ was blocked. Steps 1--2 retract stale values (underconsistent); the remaining
  steps are ordinary expansions. ``--'' means that \LpaUpdateVertex changed nothing.}
  \label{tab:ch05-lpastar-repair}
  \begin{tabular}{@{}rllll@{}}
  \toprule
  Step & Popped $u$ (key) & Case & $\gcost(u)$ after & Updated: $\rhs$, key\\
  \midrule
  %%NUM-LPA-SECOND
  \bottomrule
  \end{tabular}
\end{table}

The path is read off at the end by walking backward from $T$: from each cell move to
the predecessor that attains the minimum in \cref{eq:ch05-rhs}. In the right panel of
\cref{fig:ch05-lpastar-example} this gives $T \to (4,0) \to (3,0) \to (2,0) \to (2,1)
\to (1,1) \to S$, the blue path of cost $6$ (the symmetric path through the top row has
the same cost; the tie-break picks the bottom row).
